\documentclass[12pt]{article}

\usepackage{bbm}
\usepackage{fancyhdr}
\usepackage{titling}
\usepackage{enumitem}
\usepackage{amsmath}
\usepackage{dsfont}
\usepackage{amsfonts}
\usepackage{amssymb}
\usepackage{amsthm}
\usepackage{geometry}
\usepackage{tikz}
\usepackage{algorithm}
\usepackage{algpseudocode}
\geometry{verbose,tmargin=1in,bmargin=1in,lmargin=1in,rmargin=1in}
\renewcommand{\vec}[1]{\mathbf{#1}}
\newcommand{\dec}{\text{dec}}
\newcommand{\divi}{\mathbin{\textnormal{div}}}
\title{Homework 1}
\author{Satvik Prasad}
\date{}
\newtheorem{lemma}{Lemma}[section]
\fancypagestyle{plain}{%
  \fancyhf{}
  \fancyhead[L]{\theauthor}
  \fancyhead[R]{\thedate}
}
\pagestyle{plain}
\makeatletter
\def\@maketitle{%
  \newpage
  \null
  \vskip 1em%
  \begin{center}%
    \let \footnote \thanks
    {\LARGE \@title \par}%
    \vskip 1em%
  \end{center}%
  \par
\vskip 1em}
\makeatother
\usepackage{cmbright}

\newcommand{\mat}[1]{\mathbf{#1}}

% ---------- Probability theory ----------
% Core operators (optional argument-free; use as \E{X}, \Var{X}, \Cov{X}{Y})
\DeclareMathOperator{\EOp}{\mathbb{E}}
\DeclareMathOperator{\VarOp}{Var}
\DeclareMathOperator{\CovOp}{Cov}
\DeclareMathOperator{\CorrOp}{Corr}

\renewcommand{\P}[1]{\mathbb{P}\left(#1\right)}        % \P{A}
\newcommand{\E}[1]{\EOp\left[#1\right]}                % \E{X}
\newcommand{\Var}[1]{\VarOp\left(#1\right)}            % \Var{X}
\newcommand{\Cov}[2]{\CovOp\left(#1,\,#2\right)}       % \Cov{X}{Y}
\newcommand{\Corr}[2]{\CorrOp\left(#1,\,#2\right)}     % \Corr{X}{Y}
\newcommand{\SD}[1]{\operatorname{SD}\left(#1\right)}  % \SD{X}

% Subscripted versions, e.g. \Esub{\theta}{X}, \Psub{\theta}{A}
\newcommand{\Esub}[2]{\EOp_{#1}\left[#2\right]}
\newcommand{\Psub}[2]{\mathbb{P}_{#1}\left(#2\right)}

% Conditioning, indicators, independence
\newcommand{\given}{\,\middle|\,}                      % \E{X \given Y}
\newcommand{\ind}[1]{\mathbbm{1}\left\{#1\right\}}     % \ind{X > 0}
\newcommand{\indep}{\perp\!\!\!\perp}                  % X \indep Y
\newcommand{\iid}{\overset{\text{iid}}{\sim}}          % X_i \iid \Normal(0,1)

% Convergence
\newcommand{\convp}{\xrightarrow{\;p\;}}               % in probability
\newcommand{\convd}{\xrightarrow{\;d\;}}               % in distribution
\newcommand{\convas}{\xrightarrow{\text{a.s.}}}        % almost surely
\newcommand{\convLp}[1]{\xrightarrow{L^{#1}}}          % in L^p

% Probability spaces
\newcommand{\F}{\mathcal{F}}                           % sigma-algebra
\newcommand{\B}{\mathcal{B}}                           % Borel sigma-algebra
\newcommand{\pspace}{(\Omega, \F, \mathbb{P})}

% Common distributions
\newcommand{\Normal}{\mathcal{N}}
\newcommand{\Bern}{\operatorname{Bernoulli}}
\newcommand{\Bin}{\operatorname{Binomial}}
\newcommand{\Pois}{\operatorname{Poisson}}
\newcommand{\Unif}{\operatorname{Uniform}}
\newcommand{\Exp}{\operatorname{Exp}}
\newcommand{\Geom}{\operatorname{Geometric}}
\newcommand{\Gam}{\operatorname{Gamma}}
\newcommand{\Beta}{\operatorname{Beta}}

% Number sets
\newcommand{\R}{\mathbb{R}}
\newcommand{\N}{\mathbb{N}}
\newcommand{\Z}{\mathbb{Z}}


\begin{document}
\subsection{Green Circle}
\begin{enumerate}
  \item [1] The expression
    $$
    \sum_{k = 0}^{\lfloor n/ 2\rfloor} {n \choose 2 j}
    $$
    counts the number of even sized subsets of $[n].$ Using a
    generating function, we can see
    $$
    (1+x)^n = \sum_{k = 0}^{n} {n \choose k} x^k, \qquad (1-x)^n
    =\sum_{k = 0}^{n} {n \choose k}(-1)^k x^k.
    $$
    Since we want odd terms to cancel out, we get
    $$
    \dfrac{(1+x)^n + (1-x)^n}{2} = \sum_{j = 0}^{\lfloor n/2 \rfloor
    }{n \choose 2j} x^k,
    $$
    and letting $x = 1$ we get
    $$
    2^{n-1} = \sum_{j = 0}^{\lfloor n /2 \rfloor} {n \choose 2j}.
    $$

  \item [3] We desire to evaluate
    $$
    C_n = \sum_{k = 0}^{n}(-1)^k{n \choose k}^2 = \sum_{k = 0}^{n}
    (-1)^k {n \choose k}{n \choose n - k }
    $$
    for every $n \geq k.$

    Consider the two generating functions
    $$
    A(x) = \sum_{k = 0}^{n} (-1)^k{n \choose k}x^k, \qquad B(x) =
    \sum_{k = 0}^{n} {n \choose k} x^k.
    $$

    Then, $$[x^{n}](AB) = \sum_{k = 0}^{n} (-1)^k {n \choose k}{n
    \choose n - k} = C_n.$$

    Simplifying $A, B,$ we get
    $$
    A(x) = (1-x)^n, \qquad B(x) = (1+x)^n,
    $$
    meaning we have
    $$
    C_n = [x^n] (1-x^2)^n = [x^n] \sum_{k = 0}^{n} (-1)^k {n \choose k} x^{2k}.
    $$
    Thus, if $n \equiv 0 \pmod 2,$ then
    $$
    C_n = (-1)^{n/2} {n \choose n/2}
    $$
    otherwise $C_n = 0.$

  \item [4] We have $a_0 = 0$ and $a_n = 2a_{n -1} + 1$ for $n \geq
    1.$ Define the generating function
    \begin{align*}
      A(x) = \sum_{k \geq 0} a_k x^k &= \sum_{k \geq 1} a_kx^k \\
      &= \sum_{k \geq 1} (2a_{k - 1} + 1)x^k \\
      &=2x\sum_{k \geq 0} a_{k} x^k + \sum_{k \geq 1} x^k.
    \end{align*}
    For $|x| < 1,$ we get
    $$
    A(x)  = 2xA(x) + \dfrac{x}{1-x},
    $$
    or equivalently,
    $$
    A(x) = \dfrac{x}{(1-x)(1-2x)} = \dfrac{1}{1-2x} - \dfrac{1}{1-x},
    $$
    producing generating function
    $$
    A(x) = \sum_{k \geq 0} (2x)^k - x^k = \sum_{k \geq 0}(2^k - 1) x^k,
    $$
    yielding $a_k = 2^k -1 .$
\end{enumerate}

\subsection{Blue Square}
\begin{enumerate}
  \item [7] We look for a simple formula for
    $$
    \sum_{j = 0}^{\lfloor n/3 \rfloor} {n \choose 3j}.
    $$
    We can achieve this by summing generating functions where only
    the terms with $k \equiv 0 \pmod 3$ survive.

    Notice that if we have
    $$
    \omega^3 = 1
    $$
    for some complex $\omega,$ rearranging yields
    $$
    (\omega - 1)(\omega^2 + \omega + 1) = 0.
    $$
    Given $\omega \neq 1,$ we have
    $$
    \omega^2 + \omega + 1 = 0.
    $$

    Consider $\omega^{2k} + \omega^k + 1$ for some integer $k.$ Then,
    if $k \equiv 0 \pmod 3,$ we have $\omega^{2k} + \omega + 1 = 1 +
    1 + 1 = 3.$ Otherwise, if $k \equiv 1$ or $k \equiv 2 \pmod 3,$
    then $\{w^{2k}, w^{k}\} = \{w^2, w\},$ meaning $$
    w^{2k} + w^{k} + 1 = 0.
    $$

    Thus, using this we can show
    $$
    (1 + 1)^n + (1 + \omega)^n + (1 + \omega^2)^n = \sum_{k = 0}^{n}
    {n \choose k} (1 + \omega^k + \omega^{2k})x^k = 3\sum_{k =
    0}^{\lfloor n/3 \rfloor } {n \choose k}x^k,
    $$
    meaning
    $$
    \sum_{j = 0}^{ \lfloor n/3 \rfloor} {n \choose 3j} = \dfrac{2^n +
    (1 + e^{i\pi/3})^n + (1 + e^{-i\pi/3})^n}{3}.
    $$

    Simplifying, we have
    \begin{align*}
      e^{\frac{in\pi}{6}}\left(2\cos\frac{\pi}{6}\right)^n +
      e^{-\frac{in\pi}{6}}\left(2\cos\frac{\pi}{6}\right)^n
      &= 3^{\frac{n}{2}}\left(e^{\frac{in\pi}{6}} +
      e^{-\frac{in\pi}{6}}\right) \\
      &= 3^{\frac{n}{2}} \cdot 2\cos\left(\frac{n\pi}{6}\right).
    \end{align*}

  \item [9] We have $c_n$ counting colored compositions of sum $n.$

    Any such composition will have a last element $k$ with $k$
    different options for its color. Hence, we have
    \begin{align*}
      c_n = \sum_{k = 1}^{n} k c_{n - k}
    \end{align*}
    for all $n \geq 1.$

    Defining a generating function, we get
    $$
    C(x) = 1 + \sum_{k \geq 1} c_k x^k = 1 + \sum_{k \geq 1} \sum_{j
    = 1}^{k} jc_{k - j} x^k.
    $$
    Swapping the order of summation, we get
    \begin{align*}
      C(x) &= 1 + \sum_{j \geq 1} j\sum_{k \geq j} c_{k -j} x^k \\
      &=1 +  \sum_{j \geq 1} jx ^j\sum_{k \geq 0} c_{k} x^{k} \\
      &= 1 + C(x)\sum_{j \geq 1} jx^j \\
      &= 1+ xC(x) \sum_{j \geq 0} (j+1)x^{j} \\
      &= 1 + C(x) \dfrac{x}{(1-x)^2}
    \end{align*}
    for all $|x| < 1.$ Hence, we get
    \begin{align*}
      C(x) = \dfrac{(1-x)^2}{(1-x)^2 - x} &= 1 + \dfrac{x}{x^2 -3x + 1} \\
      &= 1  + \dfrac{1}{\sqrt{5}} \left(
        \dfrac{1}{1 - \alpha x} - \dfrac{1}{1-\beta x}
      \right),
    \end{align*}
    where $\alpha = (3 + \sqrt 5)/2,$ and $\beta = (3 - \sqrt 5)/2.$
    Expanding, we get
    $$
    C(x) = 1 + \sum_{j \geq 0} \dfrac{\alpha^j - \beta^j}{\sqrt 5} x^j,
    $$
    yielding
    $$
    c_0 = 1, \qquad c_j = \dfrac{(3 + \sqrt 5 )^j - (3 - \sqrt
    5)^j}{2^j\sqrt 5}.
    $$

    Since $\alpha^2 = 3\alpha -1 $  and $\beta^ 2 = 3\beta - 1,$ ewe get
    \begin{align*}
      c_n &= \dfrac{\alpha^{n-2}(3 \alpha - 1) - \beta^{n-2}(3 \beta
      - 1)}{\sqrt 5} \\
      &=3 \dfrac{\alpha^{n-1} - \beta^{n-1}}{\sqrt 5} -
      \dfrac{\alpha^{n-2} - \beta^{n-2}}{\sqrt 5} \\
      &= 3 c_{n-1} - c_{n-2}
    \end{align*}
    as required.
\end{enumerate}

\subsection{Black Diamond}
\begin{enumerate}
  \item [10] Write the generating function for the number of ways to
    get sum $k$ when rolling two distinct dice as
    $$
    G(x) = (x + \ldots +x ^6)^2 = x^2 (1 + \ldots + x^5)^2.
    $$
    We want to find an alternative factorization of this generating
    function. Write
    $$
    1 + \ldots + x^5 = \dfrac{x^6 - 1}{x - 1}.
    $$
    If we have $x^6 = 1,$ we get $x^6 = e^{2k\pi i},$ meaning
    $$
    x = e^{k\pi i/3},
    $$
    yielding factorization
    $$
    (x - e^{i\pi/3}) (x - e^{-i\pi/3}) (x - e^{2i\pi/3})(x -
    e^{-2i\pi/3}) (x + 1),
    $$
    or equivalently,
    $$
    (x^2 - x + 1)(x^2 + x +1)(x + 1),
    $$
    and $$
    G(x) = (x^2 - x + 1)^2 (x^2 + x+1)^2 (x+1)^2 x^2.
    $$
    Hence, to construct an alternative factorization, let us take factors
    $$
    G_1(x) = (x^2 - x+ 1)^2(x^2 + x+ 1) (x+1) x = x^8 + x^6 + x^5 +
    x^4 + x^3 + x,
    $$
    and $$
    G_2(x) = (x^2 + x+ 1)(x+1)x = x^4 + 2x^3 + 2x^2 + x,
    $$
    corresponding to the dice with faces $\{8,6, 5, 4, 3, 1\}$ and
    $\{1, 2, 2, 3, 3, 4\}.$

  \item [12] For odd primes $p,$ we want to find the number of
    subsets of $[p]$ such that their sum is divisible by $p.$ The
    term in the generating function
    $$
    [z^k]\underbrace{\prod_{j = 1}^{j = p} (1  + z^{j})}_{:=G(z)}
    $$
    will count the number of ways to get a subset with sum $k.$
    Hence, we want to find
    $$
    C_p = \sum_{k = 0}^{(p + 1) / 2} [z^{kp}] \prod_{j = 1}^{j = p} (1 + z^j).
    $$

    Notice that the root of unity $\omega^{p} = 1$ can help us cancel
    out terms that are not divisible by $p.$ Let $\omega = e^{2\pi i /p}.$ Then,
    $$
    1 + \omega^k + (\omega^k)^2 + \ldots + (\omega^k)^{p - 1} = 0.
    $$
    for all $k \not \equiv 0 \pmod p,$ and $p - 1$ otherwise. Thus,
    $$
    \sum_{j = 0}^{p-1}G(w^j) = \sum_{k = 0}^{(p + 1)/2} (p-1)[z^{kp}]
    G(z) = (p-1)C_p.
    $$

    Note that
    $$
    G(1) = \prod_{j = 1}^{p} (1+1^j) = 2^p.
    $$

    Otherwise,
    \begin{align*}
      G(\omega^r) &= \prod_{j = 1}^{p}(1 + (\omega^r)^j) \\
      &= 2 \prod_{j = 1}^{p-1} (1 + \omega^{rj}).
    \end{align*}

    Given an integer $1 \leq r \leq p-1,$ let us look at $jr \pmod p$
    for all $1 \leq j \leq p - 1.$ Assume for the sake of
    contradiction that there are two $1 \leq j_1 <  j_2 \leq p - 1$ such that
    $$
    j_1r \equiv  j_2 r \pmod p.
    $$
    Then,
    $$
    r(j_2 - j_1) \equiv  0 \pmod p,
    $$
    meaning $p$ either divides $r$ or $j_2 - j_1$ by Euclid's lemma.
    Neither is possible, since both quantities are within $[1, p -
    1].$ Thus, the image of $jr \pmod p$ must be $\{1, 2, \ldots, p - 1\}.$

    Hence, we can write
    \begin{align*}
      G(\omega^r) &= 2 \prod_{j = 1}^{p-1} (1 + \omega^{j}).
    \end{align*}
    Since
    $$
    \prod_{j = 1}^{p-1} (x-\omega^j)
    $$
    has roots $\omega, \omega^2, \ldots, \omega^{p-1},$ we can
    conclude it is equal to
    $$
    \dfrac{x^p - 1}{x- 1} = 1 + x + \ldots + x^{p-1}.
    $$
    Letting $x = -1,$ we have
    $$
    (-1)^{p-1}\prod_{j = 1}^{p-1}(x + \omega^j) = 1 -  x + \ldots +
    (-1)^{p-1}x^{p-1}.
    $$
    Since $p$ is odd, we finally get
    $$
    \prod_{j = 1}^{p-1}(1 + \omega^j) = 1,
    $$
    i.e., $G(\omega^r) = 2$ for $1 \leq r \leq p -1.$ Finally, this means
    $$
    C_p = \dfrac{2^p + (p-1) \cdot 2  }{p-1} = \dfrac{2^p}{p-1} + 2.
    $$
\end{enumerate}

\subsection{Reflection}
Had to use lots of roots of unity for this problem set, which I
learned back in high-school and forgot after that. Got
$$
\sum_{k = 1}^{p - 1} (1 + \cos(2\pi i j/p))
$$
as one of my polynomials by multiplying conjugate pairs, but realized
that the easier way to solve it was to write $\prod_{j =1}^{p-1} 1 +
\omega^j$ as a value on the polynomial
$$
\dfrac{x^p - 1}{x-1}
$$
sincee $\omega$ is a root of unity.
\end{document}
